\begin{table}[!htbp]
\centering
\footnotesize
\setlength{\tabcolsep}{3.5pt}
\caption*{\textbf{Table 6.} Guaranteed reductions in uncertainty under residence reporting}
\begin{tabularx}{\linewidth}{lCCCCC}
\toprule
Election & Current range width & Proposal: out-of-zone early ballots by neighborhood, share removed & Proposal: maximum remaining width & All pooled modes by neighborhood, share removed & All pooled modes: maximum remaining width \\
\midrule
2022 & 0.25 & 27\% & 0.18 & 57\% & 0.11 \\
2024 & 0.20 & 30\% & 0.14 & 49\% & 0.11 \\
2025 & 0.23 & 32\% & 0.16 & 61\% & 0.09 \\
2026 & 0.10 & 57\% & 0.04 & 69\% & 0.03 \\
\bottomrule
\end{tabularx}
\par\vspace{2pt}{\raggedright\footnotesize\textit{Notes.} All columns describe the width of the within-district slope range or the guaranteed fraction of that width removed. Main assumptions: Table 2, column E. Proposal: out-of-zone early ballots reported as aggregate counts by the voter's registered neighborhood, at-home and overseas ballots kept pooled. All pooled modes: out-of-zone early, at-home and overseas ballots by neighborhood. Maximum remaining width: the largest difference in the within-district slope between two allocations that the report could not tell apart, which bounds the range left whatever the true allocation. Other assumptions and headcount reporting: Online Appendix Table C12.\par}
\end{table}
